\appendix

\section{Passive selected-port Gramian bound}
\label{app:gramian}

For the passive endpoint of Eq.~\eqref{eq:passiveio}, define the complete controllability Gramian
\begin{equation}
P_{\rm all}
=\int_0^\infty\dd t\,
\ee^{At}K^\dagger K\ee^{A^\dagger t}.
\end{equation}
It solves
\begin{equation}
AP_{\rm all}+P_{\rm all}A^\dagger+K^\dagger K=0.
\end{equation}
Equation~\eqref{eq:Apassive} also gives
\begin{equation}
AI+IA^\dagger+K^\dagger K=0.
\end{equation}
For stable $A$, the Lyapunov solution is unique, hence
\begin{equation}
P_{\rm all}=I.
\end{equation}
For a selected input subset $u$,
\begin{equation}
P_u=\int_0^\infty\dd t\,
\ee^{At}K_u^\dagger K_u\ee^{A^\dagger t}.
\end{equation}
Since $0\le K_u^\dagger K_u\le K^\dagger K$, the integral representation gives
\begin{equation}
0\le P_u\le I.
\end{equation}
The standard Gramian identity then yields
\begin{equation}
\|S_{g\leftarrow u}\|_2^2
=\Tr(K_gP_uK_g^\dagger)
\le\Tr(K_g^\dagger K_g).
\end{equation}
The receiver-side statement follows analogously from the corresponding observability or dual controllability Gramian.

\section{Cumulative quadrupole spectral bound}
\label{app:sumrule}

For a passive diagonal state
\begin{equation}
\rho=\sum_m p_m|m\rangle\langle m|,
\end{equation}
with $p_m\ge p_n$ for $E_m<E_n$, the double-commutator identity for a Hermitian operator $F$ is
\begin{equation}
\frac12\Tr\rho[F,[H,F]]
=\sum_{m<n}(p_m-p_n)(E_n-E_m)|F_{mn}|^2.
\end{equation}
For the five orthonormal STF mass-quadrupole components $Q_a$ of ordinary nonrelativistic matter,
\begin{equation}
\sum_a\frac12\Tr\rho[Q_a,[H,Q_a]]
=\frac{10}{3}\hbar^2\langle I\rangle_\rho.
\end{equation}
Define the positive spectral measure
\begin{equation}
\dd\mu_Q(\omega)
=\sum_{m<n,a}
(p_m-p_n)|Q^a_{mn}|^2
\delta(\omega-\omega_{nm})\dd\omega.
\end{equation}
Then
\begin{equation}
\int_0^\infty\omega\,\dd\mu_Q(\omega)
=\frac{10}{3}\hbar\langle I\rangle_\rho.
\end{equation}
The positive gravitational transition-rate weight below $\Omega$ is
\begin{equation}
K_g(\Omega)
=\frac{2G}{5\hbar c^5}
\int_0^\Omega\omega^5\dd\mu_Q(\omega).
\end{equation}
Since $\omega^5\le\Omega^4\omega$ on the selected interval,
\begin{equation}
\boxed{
K_g(\Omega)
\le
\frac{4G}{3c^5}
\langle I\rangle_\rho\Omega^4.}
\end{equation}
For the coherent linear-network theorem we use the ground-state one-quantum coupling resource; finite-temperature population factors affect net absorption and channel noise and are not silently identified with the state-independent oscillator damping matrix.

\section{Independent aligned-plus TT integral}
\label{app:TTcheck}

Take $Q=\mathrm{diag}(1,-1,0)$ and propagation along $z$. With the usual spherical transverse basis and unit-normalized polarization tensors,
\begin{align}
Q:\epsilon_+
&=\frac{1+\cos^2\theta}{\sqrt2}\cos2\phi,\\
Q:\epsilon_\times
&=-\sqrt2\cos\theta\sin2\phi.
\end{align}
For the normalized one-graviton angular mode in Eq.~\eqref{eq:normalizedmode}, azimuthal integration gives
\begin{equation}
\sum_\lambda\int_0^{2\pi}\dd\phi\,
|u_{Q,\lambda}|^2
=
\frac5{32}(1+6x^2+x^4),
\qquad x=\cos\theta.
\end{equation}
Thus the translated overlap is
\begin{equation}
S(z)
=\frac5{32}\int_{-1}^{1}
(1+6x^2+x^4)\ee^{izx}\dd x.
\end{equation}
The endpoint weight is $5/4$ at both $x=\pm1$. A single integration by parts therefore gives the outgoing contribution
\begin{equation}
S_+(z)
=-\frac{5i}{4z}\ee^{iz}+O(z^{-2}).
\end{equation}
Evaluating the polynomial Fourier integral exactly gives
\begin{equation}
S_+(z)
=-\frac{5i}{4}\frac{P(z)\ee^{iz}}{z^5},
\end{equation}
where
\begin{equation}
P(z)=3-3iz-3z^2+2iz^3+z^4.
\end{equation}
Consequently,
\begin{equation}
|S_+(z)|^2
=
\frac{25}{16z^2}
\left(
1-\frac2{z^2}+\frac3{z^4}
-\frac9{z^6}+\frac9{z^8}
\right).
\end{equation}
This check fixes the leading $25/16$ coefficient directly from normalized angular modes.

\section{Exact two-pole integral and symmetric optimum}
\label{app:twopole}

Substitution of $a=\kappa_A/2$ and $b=\kappa_B/2$ into
\begin{equation}
\int_{-\infty}^{\infty}\frac{\dd\Omega}
{(\Omega^2+a^2)(\Omega^2+b^2)}
=\frac{\pi}{ab(a+b)}
\end{equation}
gives Eq.~\eqref{eq:exactEBP} immediately. To prove Eq.~\eqref{eq:twopolecut}, note first that
\begin{equation}
4\kappa_{g,B}\kappa_{\rm out}
\le(\kappa_{g,B}+\kappa_{\rm out})^2
\le\kappa_B^2
\end{equation}
and $\kappa_{\rm in}\le\kappa_A$. Hence
\begin{equation}
\Gamma_{\rm EBP}
\le
\eta_{\rm prop}\kappa_{g,A}
\frac{\kappa_B}{\kappa_A+\kappa_B}
\le
\eta_{\rm prop}\kappa_{g,A}.
\end{equation}
Exchanging $A$ and $B$ gives the receiver-side bound.

For the symmetric lossless family, define
\begin{equation}
f(x)=\frac{2x^2}{(1+x)^3}.
\end{equation}
Then
\begin{equation}
f'(x)=\frac{2x(2-x)}{(1+x)^4},
\end{equation}
so the unique positive maximum occurs at $x=2$, with $f(2)=8/27$.
